\chapter{Memorie Cache}
\label{ch:memorie-cache}

\section{Introduzione e Motivazione}
Le memorie cache sono memorie di piccole dimensioni ma ad altissima velocità, interposte tra il processore e la memoria principale.
La necessità fondamentale di introdurre questa gerarchia deriva dall'enorme divario prestazionale tra il processore e la memoria. Attualmente, le prestazioni dei processori rispetto a quelle della memoria principale differiscono di un fattore di circa 1000. In termini pratici, nel tempo in cui il processore è in grado di eseguire migliaia di operazioni, la memoria principale è in grado di fornire una sola istruzione. Senza un sistema per mantenere dati e codice pronti e accessibili velocemente, il processore passerebbe la maggior parte del tempo in stallo, vanificando qualsiasi ottimizzazione architetturale, inclusa la pipeline. 

Il costo energetico e temporale del calcolo odierno risiede prevalentemente nel trasferimento delle informazioni: elaborare un dato, per quanto l'operazione sia complessa, costa in termini di potenza circa mille volte meno rispetto al trasferire quel dato dalla memoria al processore.

\subsection{Architettura di Von Neumann vs Harvard}
Per comprendere l'integrazione della memoria nel sistema, è fondamentale richiamare la differenza tra le architetture di base:
\begin{itemize}
    \item \textbf{Architettura Von Neumann}: Dati e istruzioni (codice, sub-routine, funzioni) risiedono in una singola memoria unificata. Il processore accede alla stessa memoria per prelevare sia il codice da eseguire sia i dati da elaborare.
    \item \textbf{Architettura Harvard}: Prevede memorie fisicamente separate per i dati e per il codice. Il processore può accedere simultaneamente alle due memorie attraverso bus distinti.
\end{itemize}
Nei sistemi moderni e specialmente nei processori RISC (come RISC-V e ARM) che implementano una pipeline, l'architettura Harvard è essenziale al livello della cache (Cache L1 Dati e Cache L1 Istruzioni separate). Ciò permette di eseguire contemporaneamente il fetch di un'istruzione e la lettura/scrittura di un dato, evitando conflitti strutturali che costringerebbero la pipeline a stallare.

\section{Gerarchia delle Memorie e Principio di Località}
Tra la memoria di immagazzinamento di massa (Hard Disk/SSD) e il processore, viene interposta una gerarchia di memorie per garantire un accesso rapido alle informazioni. 

Analizzando il sistema partendo dall'interno verso l'esterno, troviamo:
\begin{enumerate}
    \item \textbf{Registri della CPU}: Elementi di memoria interni (flip-flop) a cui la CPU accede alla massima frequenza di sistema (centinaia di picosecondi). La loro capacità è molto limitata, tipicamente intorno a un kilobyte (es. 32 registri da 32 bit più i registri di controllo).
    \item \textbf{Cache L1 (Livello 1)}: Memoria molto veloce (accesso in 1-2 nanosecondi) e di dimensioni ridotte (es. 64 KB). Tipicamente divisa in Dati e Istruzioni.
    \item \textbf{Cache L2 (Livello 2) e L3 (Livello 3)}: Memorie progressivamente più grandi (da centinaia di KB a diversi MB) ma via via più lente.
    \item \textbf{Memoria Principale (RAM)}: Memoria di lavoro capiente (diversi GB), ma con tempi di accesso lenti, nell'ordine dei 50-100 nanosecondi.
\end{enumerate}

% FIGURA ORIGINALE DA RECUPERARE DALLE SLIDE
% Slide di provenienza: 4
% Tipo: diagramma
% Descrizione: Diagramma dei livelli di memoria per dispositivi mobili, laptop e server, con relativi tempi di accesso e capacità.
% Punto di inserimento: In questo punto, a supporto della classificazione dei livelli di memoria.
% File da inserire: fig_memory_levels

L'assunzione alla base dell'architettura gerarchica è il \textbf{principio di inclusività}: tutte le informazioni contenute nei registri dovrebbero trovarsi nella cache L1; tutte le informazioni della cache L1 devono essere contenute nella cache L2, e così via fino alla memoria principale. Se un dato non è in L1 (situazione di \textit{miss}), il sistema lo cerca in L2, poi in L3, con un progressivo e drastico calo delle prestazioni a ogni salto.

\begin{remark}[Inizializzazione nei Simulatori]
Quando si analizzano le prestazioni tramite un simulatore, le prime fasi di avvio (es. la routine \texttt{\_start}) comportano numerosi \textit{cache miss} poiché la cache è inizialmente vuota. Per valutare il tempo reale di un algoritmo, è opportuno sottrarre i colpi di clock dell'inizializzazione (circa 208 cicli in alcuni simulatori) o assicurarsi di iniziare il conteggio dal \texttt{main}, operando una lettura preliminare dei dati per caricarli in cache.
\end{remark}

\subsection{Principio di Località}
L'efficacia della cache si basa sulla località dei riferimenti osservata nell'esecuzione della maggior parte dei programmi. Si manifesta in due forme:
\begin{itemize}
    \item \textbf{Località Temporale}: Se al tempo $t$ un programma accede a una determinata cella di memoria, è altamente probabile che acceda nuovamente alla stessa cella in un istante prossimo $t_0 + \Delta t$. Per questo motivo, quando un dato viene letto, esso viene mantenuto in cache.
    \item \textbf{Località Spaziale}: Se al tempo $t_0$ il programma accede all'indirizzo $x$, è molto probabile che a breve accederà agli indirizzi contigui $x \pm e$ (ad esempio, gli elementi successivi di un vettore). Per sfruttare questa proprietà, quando si verifica un \textit{miss}, la cache non carica dalla memoria soltanto il byte richiesto, ma l'intera \textbf{linea di cache} (es. 32 o 64 byte contigui).
\end{itemize}

\subsection{Prestazioni e Hit Ratio}
Per quantificare le prestazioni del sistema di memoria, si definiscono i seguenti parametri:
\begin{itemize}
    \item $h$: \textbf{Hit ratio} (tasso di successo), ovvero la percentuale di volte in cui il dato richiesto viene trovato in cache.
    \item $C$: Tempo di accesso alla cache.
    \item $M$: Tempo di accesso alla memoria principale (in caso di \textit{miss}).
\end{itemize}

Il tempo medio di accesso alla memoria ($t_{\text{ave}}$) si calcola come:
\[t_{\text{ave}} = h \cdot C + (1 - h) \cdot M\]

Valori normali di $h$ si attestano sull'80\%-90\% ($0.8 - 0.9$).
Se, per esempio, abbiamo un hit ratio dell'80\% ($h=0.8$), un tempo di accesso alla cache di 2 ns ($2 \times 10^{-9}$ s) e un tempo di accesso alla RAM ipotizzato, per enfatizzare il ritardo, a $50 \mu$s, il calcolo sarà:
\[T_{\text{avg}} = 0.8 \times (2 \cdot 10^{-9}) + 0.2 \times (50 \cdot 10^{-6}) \approx 10.001 \mu s\]
Si nota come anche una piccola frazione di miss ($0.2$) domini completamente il tempo medio di accesso a causa dell'enorme scarto temporale tra $C$ e $M$.

\section{Organizzazione Interna della Cache}
Come quasi tutti i dispositivi digitali, la cache è divisa in una parte dedicata ai dati e una parte di controllo. Più specificamente, una memoria cache è strutturata in \textbf{linee} (o blocchi) ed è costituita da due componenti principali:
\begin{enumerate}
    \item \textbf{Array dei Dati (Data Array)}: Contiene i blocchi di memoria effettivi caricati dalla memoria principale. Ogni linea contiene una quantità $B$ di byte (es. 4, 32 o 128 byte).
    \item \textbf{Array di Direttorio (Tag Array)}: Contiene le informazioni necessarie per identificare univocamente da quale indirizzo della memoria principale provengono i dati salvati nella linea corrispondente.
\end{enumerate}

% FIGURA ORIGINALE DA RECUPERARE DALLE SLIDE
% Slide di provenienza: 9 e 10
% Tipo: diagramma
% Descrizione: Diagramma dell'organizzazione interna della cache con array dei dati, array di direttorio, tag, e bit di validità.
% Punto di inserimento: Qui, per illustrare visivamente la divisione in linee.
% File da inserire: fig_cache_org

Oltre ai dati e al tag, ogni linea di cache contiene alcuni bit di controllo fondamentali:
\begin{itemize}
    \item \textbf{Bit di Validità (V - Validity bit)}: Indica se il dato presente nella linea è valido. All'accensione del sistema (power-on) o in caso di reset, la memoria cache contiene valori casuali (dati "sporchi"). Per evitare che la CPU legga informazioni scorrette, il bit V viene inizializzato a 0. Quando un blocco di memoria valido viene caricato nella linea, V viene settato a 1.
    \item \textbf{Dirty bit (D)}: Ricorda se il blocco caricato in cache è stato modificato rispetto alla copia presente nella memoria principale (utilizzato nelle politiche di \textit{write-back}). Se D = 1, i byte in cache sono più recenti e dovranno essere riscritti in memoria principale prima che la linea venga sovrascritta.
\end{itemize}

\subsection{Ricerca di un blocco dati nella Cache}
Per comprendere se un dato è presente in cache, la logica di controllo riceve l'indirizzo a 32 bit generato dalla CPU (indirizzo di memoria) e lo suddivide in tre campi: \textbf{Tag}, \textbf{Indice (Index)} e \textbf{Offset}.

% FIGURA ORIGINALE DA RECUPERARE DALLE SLIDE
% Slide di provenienza: 15
% Tipo: schema
% Descrizione: Schema della scomposizione dell'indirizzo della CPU in Tag, Index e Offset.
% Punto di inserimento: Qui.
% File da inserire: fig_cpu_address

\begin{enumerate}
    \item \textbf{Offset}: Sono i bit meno significativi dell'indirizzo. Servono per individuare il byte specifico all'interno della linea di cache. Se una linea è ampia $2^n$ byte, l'offset sarà composto da $n$ bit. Questo campo non viene salvato nel direttorio, poiché indica solo una posizione relativa all'interno del blocco dati già caricato.
    \item \textbf{Indice (Index)}: Determina a quale linea (o insieme di linee, detto \textit{set}) della cache accedere. Se la cache possiede $2^m$ insiemi, l'indice utilizzerà $m$ bit.
    \item \textbf{Tag}: Costituisce la parte più significativa dell'indirizzo ed è l'effettiva informazione che viene salvata nell'array di direttorio. La sua dimensione è data dai restanti bit: $32 - (n + m)$. 
\end{enumerate}

Il procedimento di ricerca avviene in questo modo:
\begin{enumerate}
    \item La cache usa l'Indice per selezionare la riga (set) corrispondente.
    \item Controlla il Bit di Validità di quella riga. Se $V=0$, si ha immediatamente un \textit{miss}.
    \item Se $V=1$, confronta il Tag memorizzato nel direttorio della riga con il campo Tag dell'indirizzo richiesto.
    \item Se corrispondono, si ha un \textbf{Hit} (e l'offset estrae il byte richiesto); se non corrispondono, si ha un \textbf{Miss}.
\end{enumerate}

\begin{remark}[Dimensione Reale della Cache]
Quando un costruttore vende una cache, ad esempio, da 64 KB, si riferisce unicamente alla capacità dell'Array dei Dati. La dimensione fisica reale implementata sul chip è molto più grande, poiché a ogni linea va sommato lo spazio per salvare il Tag, il bit V e il bit D, moltiplicato per il numero totale di linee.
\end{remark}

\subsection{Posizionamento della Cache nel Sistema}
La cache viene collocata logicamente tra la CPU e il bus di sistema, piuttosto che tra il bus e la memoria principale.
Questo posizionamento strategico comporta che, in caso di \textit{hit}, il processore legga o scriva il dato senza alcun transito di informazioni sul bus di sistema. Ciò riduce drasticamente il carico sul bus, lasciandolo libero per periferiche e DMA, e risulta indispensabile nelle architetture multiprocessore per permettere ad ogni core di lavorare indipendentemente sui propri dati.

\section{Tecniche di Mappatura (Mapping)}
Il meccanismo di mapping definisce in quale linea della cache un determinato blocco di memoria principale deve essere posizionato. Questo processo deve essere rapido e richiedere un hardware a basso costo.

\subsection{Mappatura Diretta (Direct Mapping)}
Nella mappatura diretta, ogni blocco di memoria $i$ viene associato staticamente e univocamente a una linea (o set) $k$ della cache mediante l'operazione di modulo:
\[k = i \bmod N\]
dove $N$ è il numero totale di linee della cache.

In termini hardware, questa operazione non richiede una reale divisione, poiché il calcolo corrisponde semplicemente al prelievo dei bit meno significativi (l'Indice) dell'indirizzo del blocco.

\textbf{Vantaggi:} Il meccanismo è facilissimo da implementare in hardware. \\
\textbf{Svantaggi (Problema del Thrashing):} Se il programma accede frequentemente a due indirizzi di memoria che, pur differendo nei bit più alti, presentano gli stessi bit di indice (es. l'indirizzo $4$ e l'indirizzo $12$ modulo $8$ restituiscono entrambi la posizione $4$), essi competeranno per la stessa unica riga della cache. Ad ogni lettura ci sarà un \textit{miss} che obbligherà il sistema a sfrattare il dato precedente, abbattendo le prestazioni.

\subsection{Mappatura Associativa agli Insiemi (Set Associative)}
Per arginare il problema della competizione per le righe, la cache viene divisa in un numero $S$ di \textbf{insiemi (Set)}, ognuno dei quali contiene un numero $W$ di linee detto \textbf{vie (Ways)}. 
Si definisce quindi l'equazione:
\[S = \frac{N}{W}\]
dove $N$ è il numero totale di linee. Un blocco di memoria $i$ viene ora associato al set $k$:
\[k = i \bmod S\]
Una volta calcolato il set $k$, il blocco di memoria può essere salvato in \textit{qualsiasi} delle $W$ linee disponibili in quel set. Se abbiamo una cache a 2 vie ($W=2$), l'indirizzo $4$ e l'indirizzo $12$ calcoleranno lo stesso indice, ma potranno convivere tranquillamente occupando le due linee del medesimo set.

% FIGURA ORIGINALE DA RECUPERARE DALLE SLIDE
% Slide di provenienza: 28 e 29
% Tipo: diagramma
% Descrizione: Diagrammi rappresentativi del Mapping Associativo agli Insiemi e relativo esempio di allocazione di blocchi multipli nello stesso set.
% Punto di inserimento: Qui.
% File da inserire: fig_set_associative

\textbf{Effetti dell'Associatività sull'Indirizzo:}
Aumentare l'associatività a parità di dimensione della cache (Array Dati fisso) comporta che il numero di Set ($S$) diminuisca.
Di conseguenza:
\begin{itemize}
    \item La dimensione del campo \textbf{Indice} diminuisce (perché ci sono meno set da indicizzare).
    \item La dimensione del campo \textbf{Tag} aumenta (per compensare i bit persi dall'indice).
    \item Poiché il campo Tag si allunga e va salvato per ciascuna delle $N$ linee (le linee restano invariate in numero assoluto), \textbf{la dimensione fisica reale della memoria di controllo aumenta}.
\end{itemize}

\subsection{Mappatura Completamente Associativa (Fully Associative)}
In questo caso, l'Indice viene eliminato del tutto. Un blocco di memoria può risiedere in \textit{qualsiasi} linea della cache. 
Questo schema offre la massima flessibilità e minimizza gli scarti forzati (no conflitti di set), ma richiede un hardware estremamente complesso, poiché al momento della ricerca la CPU deve confrontare simultaneamente il Tag con \textit{tutte} le linee della cache (memoria CAM - Content Addressable Memory). Il campo Tag in questo caso copre quasi tutto l'indirizzo, escluso l'Offset.

\section{Algoritmi di Rimpiazzamento}
Quando si ha un \textit{miss} e tutte le vie del Set calcolato sono già occupate da dati validi, è necessario scegliere quale blocco sfrattare per fare spazio al nuovo dato. Le strategie principali sono:
\begin{itemize}
    \item \textbf{LRU (Least Recently Used)}: Rimpiazza la linea usata meno di recente. Richiede di memorizzare la "storia" degli accessi (es. tramite contatori), risultando costoso in termini di bit addizionali e logica.
    \item \textbf{FIFO (First-In First-Out)}: Rimuove semplicemente il dato entrato per primo (il più vecchio), a prescindere da quanto frequentemente venga usato. 
    \item \textbf{LFU (Least Frequently Used)}: Teoreticamente il più efficace, guarda alla frequenza di utilizzo assoluta, ma molto complesso da implementare.
    \item \textbf{Random}: Semplice ed efficace, sostituzione casuale.
\end{itemize}

\subsection{Approccio FIFO con Seconda Chance (Clock Bit)}
Per mitigare il problema del FIFO puro, che potrebbe espellere un dato molto vecchio ma costantemente utilizzato, si adotta spesso una politica con "Seconda Opportunità". Si associa a ogni linea un \textbf{Clock bit} che viene settato a $1$ qualora vi sia un "hit" su quella riga (ovvero il dato viene riutilizzato).
Quando il sistema deve espellere una riga, un puntatore circolare esamina le vie in ordine d'ingresso:
\begin{itemize}
    \item Se il Clock bit è $1$, il sistema lo porta a $0$, non espelle il dato (dandogli una seconda chance), e fa avanzare il puntatore.
    \item Se il Clock bit è $0$, il dato viene sfrattato.
\end{itemize}

\subsection{Algoritmo pseudo-LRU (pLRU)}
Poiché un vero algoritmo LRU richiederebbe, per un'architettura a 4 vie, 8 bit di storia per ogni Set (diventando troppo costoso in hardware), si adotta una sua approssimazione efficiente basata su un albero binario decisionale, aggiungendo soli 3 bit di storia ($a_0, b_0, b_1$) per set (nel caso di 4 vie).

% FIGURA ORIGINALE DA RECUPERARE DALLE SLIDE
% Slide di provenienza: 36 (fino a 45 per i vari passaggi A-E)
% Tipo: diagramma
% Descrizione: Funzionamento dell'algoritmo pLRU passo-passo con l'inserimento degli elementi A, B, C, D, E nell'albero binario.
% Punto di inserimento: Qui.
% File da inserire: fig_plru_tree

I bit indicano il percorso nell'albero per trovare la "foglia" (via) usata meno di recente. Quando una via viene scritta o letta, i bit sul percorso che conducono a essa vengono invertiti, indicando di "non andare più da quella parte" per il prossimo rimpiazzo.
Ad esempio, se i 3 bit partono da stato `0 0 0`, e vengono generati miss consecutivi (A, B, C, D):
\begin{enumerate}
    \item Arriva il dato \textbf{A}: l'albero suggerisce la via 0. Il dato viene allocato e i nodi attraversati invertono stato.
    \item Arriva il dato \textbf{B}: i nodi ora puntano alla via 1.
    \item Arriva \textbf{C}: le inversioni precedenti fanno deviare la radice verso l'altro ramo, puntando alla via 2.
    \item Arriva \textbf{D}: viene indicata la via 3.
\end{enumerate}
È molto più compatto, anche se in certi pattern anomali di lettura potrebbe finire col preservare elementi non necessari.

\section{Gestione delle Scritture e Coerenza}

\subsection{Aggiornamento della Memoria Principale}
Quando il processore scrive un dato in cache, sorge il problema di quando aggiornare la memoria principale per allinearla al nuovo valore.
Le due soluzioni sono:
\begin{itemize}
    \item \textbf{Write-Back}: Il dato viene aggiornato in memoria principale \textit{esclusivamente} nel momento in cui la linea corrispondente deve essere sfrattata dalla cache. Per ricordare che il dato è stato alterato rispetto alla RAM, la linea imposta a $1$ il suo \textbf{Dirty bit}. \\
    \textbf{Pro}: Riduce enormemente gli accessi e l'occupazione del bus.\\
    \textbf{Contro}: Sfrattare un dato costa tempo aggiuntivo (deve essere fisicamente riscritto in RAM). In sistemi multiprocessore, genera problemi di coerenza.
    \item \textbf{Write-Through}: Ogni singola operazione di scrittura viene propagata \textit{immediatamente} sia sulla cache che sulla memoria principale. \\
    \textbf{Pro}: La memoria è sempre aggiornata e coerente.\\
    \textbf{Contro}: Abbassa le prestazioni impegnando il bus, anche se limitato dal fatto che le letture statisticamente superano di molto le scritture.
\end{itemize}

Inoltre, se la scrittura del processore avviene su un indirizzo non presente in cache (\textbf{write miss}), si può decidere se:
\begin{itemize}
    \item \textbf{Write-allocate}: Caricare l'intero blocco dalla memoria alla cache, modificarlo in cache e mantenerlo lì. Usato tipicamente col \textit{Write-back}.
    \item \textbf{No write-allocate}: Inviare il dato direttamente in memoria senza caricarlo in cache. Usato solitamente col \textit{Write-through}.
\end{itemize}

\subsection{Coerenza della Memoria (Multiprocessore) e Protocollo MESI}
In architetture HPC (High Performance Computing) e nei classici sistemi cluster o multicore che condividono la memoria centrale (e una o più cache di basso livello, es. L3 condivisa tra processori aventi L1 e L2 private), il \textit{Write-back} fa sì che il core $A$ abbia un dato aggiornato nella sua cache L1 privata che il core $B$ non vede, causando letture corrotte (incoerenza).

Per mantenere la coerenza, ogni linea di cache, in aggiunta ai bit V e D, traccia il suo stato di condivisione. Il protocollo più celebre è il \textbf{MESI}:
\begin{itemize}
    \item \textbf{M - Modified}: La linea è stata modificata dal core locale ed è posseduta in modo esclusivo (il dato in RAM non è aggiornato).
    \item \textbf{E - Exclusive}: La linea è coerente con la RAM ed è presente esclusivamente in questa singola cache locale.
    \item \textbf{S - Shared}: La linea è stata letta ma non modificata; anche altre cache la possiedono e i valori combaciano tutti con la RAM.
    \item \textbf{I - Invalid}: Se il core $A$ scrive su un blocco che in precedenza era condiviso (Shared), invia sul bus un segnale che forza il blocco corrispondente in tutte le altre cache a marcare la linea come Invalida. Se il core $B$ chiederà il dato, dovrà prima costringere il core $A$ ad aggiornare la memoria o passargli direttamente il nuovo valore.
\end{itemize}

\section{Esempi Calcolati e Strutturazione Fisica}

In questa sezione si illustrano esplicitamente i calcoli atti a dimensionare tag, indice e offset, necessari per risalire alla reale grandezza fisica di una memoria cache in aggiunta all'array dei dati.

\subsection{Esempio 1: Cache Direct Mapped}
Assumiamo le seguenti caratteristiche:
\begin{itemize}
    \item Capacità Array Dati: 64 kB
    \item Mappatura: Direct Mapping
    \item Dimensione linea/blocco: 4 Byte 
    \item Indirizzi su 32 bit
\end{itemize}

\textbf{Procedimento:}
\begin{enumerate}
    \item \textbf{Calcolo dell'Offset}: Poiché ogni linea contiene 4 Byte ($2^2$), servono \textbf{2 bit} di Offset per indirizzare i singoli byte all'interno del blocco.
    \item \textbf{Calcolo dell'Indice}: Il numero di linee totali è la dimensione totale diviso la capacità della singola linea:
    \[ \frac{64 \text{ kB}}{4 \text{ Byte}} = \frac{2^6 \cdot 2^{10}}{2^2} = 2^{14} \text{ linee (16k linee)} \]
    L'Indice richiederà quindi \textbf{14 bit}.
    \item \textbf{Calcolo del Tag}: Si sottraggono i bit di Indice e Offset dall'indirizzo:
    \[ 32 - 14 - 2 = \textbf{16 bit} \text{ per il Tag} \]
\end{enumerate}

La dimensione reale fisica di questa cache non è 64 kB, bensì bisogna sommare ad ogni blocco (linea) di 4 Byte (32 bit) anche la dimensione del Tag (16 bit), tralasciando per semplicità i bit V e D:
\[ \text{Capacità Totale} = 2^{14} \text{ linee} \times (32 \text{ bit dati} + 16 \text{ bit tag}) = 2^{14} \times 48 \text{ bit} = 768 \text{ kbit} \]
Dividendo per 8, otteniamo \textbf{96 kByte}, un valore sensibilmente superiore al dato nominale (64 kB).

\subsection{Esempio 2: Cache Set-Associative in un Layout di Memoria Misto}
Consideriamo il seguente scenario, tipico di sistemi embedded, in cui il processore richiede il calcolo dei parametri di una cache con queste specifiche:
\begin{itemize}
    \item Capacità Array Dati: 4 KB
    \item Mappatura: 2-vie Set Associative (Associatività = 2)
    \item Lettura da memoria: 32 Byte per operazione. La linea di cache conterrà quindi blocchi di 32 Byte.
    \item Indirizzi su 32 bit
\end{itemize}

% FIGURA ORIGINALE DA RECUPERARE DALLE SLIDE
% Slide di provenienza: 61 a 73 (Animazione layout di memoria, codice Assembly e passaggi calcolo cache)
% Tipo: codice ed evoluzione schema
% Descrizione: Rappresentazione dello spazio di indirizzamento frammentato (Flash, vuoto, RAM) associato al calcolo passo passo di Offset, Index e Tag.
% Punto di inserimento: Qui.
% File da inserire: fig_example2_memory_layout

\textbf{Procedimento:}
\begin{enumerate}
    \item \textbf{Offset}: Per contenere 32 Byte, l'Offset necessita di \textbf{5 bit} ($2^5 = 32$).
    \item \textbf{Linee totali e Indice}: Calcoliamo il numero totale di linee:
    \[ \frac{4 \text{ KB}}{32 \text{ Byte}} = \frac{4096}{32} = 128 \text{ linee totali} \]
    Siccome la cache è a 2 vie, le linee vengono raggruppate in insiemi (Set) da due.
    \[ \text{Numero di Set} = \frac{128 \text{ linee}}{2} = 64 \text{ Set} \]
    Per indirizzare 64 Set ($2^6$) servono \textbf{6 bit} di Indice.
    \item \textbf{Tag}: Sottraiamo Indice e Offset:
    \[ 32 - 5 - 6 = \textbf{21 bit} \text{ per il Tag} \]
\end{enumerate}
La memoria extra aggiunta per il salvataggio dei Tag sarà pari a:
\[ 128 \text{ linee (totali)} \times 21 \text{ bit} = 2688 \text{ bit} = 336 \text{ Byte} \]
Sommando la logica supplementare a una cache già molto piccola come in questo esempio, l'incremento di percentuale dello spazio richiesto è piuttosto gestibile, sebbene sia richiesto un bit V per ognuna delle 128 linee.

\subsection{Esempio 3: Cache Associativa a 4-vie}
Parametri:
\begin{itemize}
    \item Dati: 4 KB
    \item Mappatura: 4-vie Set Associative
    \item Linea: 32 Byte
    \item Indirizzo: 32 bit
\end{itemize}

\textbf{Procedimento:}
\begin{enumerate}
    \item \textbf{Offset}: $2^5 = 32$ $\rightarrow$ \textbf{5 bit}.
    \item \textbf{Linee totali}: $4\text{KB}/32 = 128$ linee totali.
    \item \textbf{Indice}: 128 linee / 4 vie = 32 Set ($2^5$). Indice = \textbf{5 bit}.
    \item \textbf{Tag}: $32 - (5 + 5) = \textbf{22 bit}$.
\end{enumerate}

Calcolo dell'area logica di controllo (Direttorio):
\begin{itemize}
    \item Tag Bits totali: $128 \times 22 = 2816$ bit.
    \item Bit di Validità (V): $128$ bit.
    \item Dirty bits (D): $128$ bit.
    \item Bits per il pLRU Replacement: l'algoritmo pseudo-LRU per 4 vie richiede 3 bit per Set. Pertanto i 32 Set avrebbero $32 \times 3 = 96$ bit. \\
    \textit{Nota: nelle slide, un'approssimazione indica prudenzialmente 2 bit per linea, ottenendo $128 \times 2 = 256$ bit. Adotteremo per coerenza il risultato didattico fornito dal professore nelle slide per il calcolo finale.}
\end{itemize}

\textbf{Dimensione totale memoria cache}:
\[ 4 \text{ KB} + \text{Bit di Controllo} \]
\[ 32768 \text{ bit (Dati)} + 2816 (\text{Tag}) + 128 (\text{V}) + 128 (\text{D}) + 256 (\text{pLRU}) = 35584 \text{ bit} = 4448 \text{ Byte} \]

\vspace{1cm}
\hrule
\vspace{0.5cm}

\section*{Immagini da recuperare}
\begin{itemize}
    \item \textbf{Slide 4 - Diagramma dei livelli di memoria}: Grafico a blocchi rappresentante gerarchia per Mobile, Laptop e Server. Da inserire nella Sezione "Gerarchia delle Memorie".
    \item \textbf{Slide 9 e 10 - Organizzazione Interna}: Schema dei blocchi Data Array e Tag/Directory Array, incluse V e D flag. Da inserire nella Sezione "Organizzazione Interna della Cache".
    \item \textbf{Slide 15 - Scomposizione CPU Address}: Rettangolo raffigurante i 32 bit suddivisi in Tag, Index e Offset. Da posizionare in "Ricerca di un blocco dati nella Cache".
    \item \textbf{Slide 28 e 29 - Set Associative Mapping}: Mappatura dei blocchi di memoria su cache suddivisa a insiemi e vie, mostrato con frecce illustrative. Da inserire sotto "Mappatura Associativa agli Insiemi".
    \item \textbf{Slide 36 a 45 - Pseudo-LRU (pLRU)}: Evoluzione temporale (sequenza A-E) di un albero binario a 3 nodi decisionali. Da inserire sotto "Algoritmo pseudo-LRU".
    \item \textbf{Slide 61 a 73 - Schema layout di Memoria Misto}: Sequenza animata/strutturata che illustra il calcolo dei bit per tag/index e layout della memoria segmentata associato al codice assembly. Da inserire in "Esempio 2".
\end{itemize}

\section*{Problemi da verificare}
\begin{itemize}
    \item \textbf{Calcolo dei bit di storia per pLRU (Esempio 3)}: Nelle slide (Slide 74) è riportato il dato ``LRU Replacement = 2-bits per line = 256 bits''. L'algoritmo pLRU ad albero binario spiegato nella Sezione 5.2, implementato per una cache a 4 vie, necessiterebbe di 3 bit totali per ciascuno dei 32 set (quindi 96 bit e non 256). Si consiglia di verificare col docente se la dicitura ``2-bits per line'' fa riferimento ad un algoritmo di history replacement diverso dall'albero binario, o se si tratti di un arrotondamento cautelativo generico per il design.
\end{itemize}